% ============================================================
% DEMO — 中国市场，模块式中文简历（基本信息网格 + 色块标题）。
% 与 australia-sample 同一批事实。虚构人物 李乔丹 / Jordan Lee。
% COMPILE WITH XeLaTeX.
% ============================================================
\documentclass[10pt]{article}
\usepackage[UTF8, fontset=windows]{ctex}

\usepackage[top=0.55in,bottom=0.55in,left=0.75in,right=0.75in]{geometry}
\usepackage{hyperref}
\usepackage{enumitem}
\usepackage{tabularx}
\usepackage{setspace}
\usepackage{graphicx}
\usepackage{xcolor}

\definecolor{accent}{HTML}{215D6E}
\hypersetup{colorlinks=true, urlcolor=accent, linkcolor=accent}
\setlength{\parindent}{0pt}
\setstretch{1.02}
\pagenumbering{gobble}
\setlist[itemize]{leftmargin=1.3em, labelsep=0.5em, itemsep=1.5pt, topsep=2pt, before=\vspace{-2pt}}
\renewcommand{\labelitemi}{\textcolor{accent}{\footnotesize$\bullet$}}

\newcommand{\cnsection}[1]{%
  \vspace{9pt}%
  {\color{accent}\rule[-0.15ex]{3.5pt}{1.05em}}\hspace{7pt}{\large\bfseries\color{accent}#1}\\[1pt]%
  {\color{accent!30}\rule{\textwidth}{0.9pt}}\\[3pt]%
}
\newcommand{\cnentry}[3]{\textbf{#1}\;\;{\small\color{gray}#3}\hfill{\small\itshape #2}\\[1pt]}

\begin{document}
\vspace*{0pt}

\noindent
\begin{minipage}[b]{0.72\textwidth}
  {\fontsize{24pt}{28pt}\selectfont\bfseries\color{accent}李乔丹}\;\;{\normalsize\color{gray}（Jordan Lee）}\\[6pt]
  \setlength{\tabcolsep}{0pt}\renewcommand{\arraystretch}{1.35}
  \begin{tabularx}{\linewidth}{@{}p{0.52\linewidth}X@{}}
    求职意向：数据科学家 / 机器学习工程师 & 电话：151 0000 0000\\
    所在城市：上海                        & 邮箱：\href{mailto:jordan.lee@example.com}{jordan.lee@example.com}\\
    数据科学硕士 · 2022 届                 & 主页：\href{https://www.linkedin.com/in/example}{linkedin.com/in/example}\\
  \end{tabularx}
\end{minipage}\hfill
\begin{minipage}[b]{0.24\textwidth}\raggedleft
  \fbox{\parbox[c][3.3cm][c]{2.6cm}{\centering 一寸\\照片}}
\end{minipage}

\cnsection{个人简介}
数据科学从业者，具备应用型机器学习系统的落地经验，覆盖数据管道、模型评估到面向非技术团队的结果沟通，寻求数据科学 / 机器学习方向岗位。

\cnsection{工作经历}
\cnentry{Example Analytics（墨尔本）}{2023.02 - 至今}{数据科学家}
\begin{itemize}
  \item 搭建结合统计方法与机器学习的异常检测管道，将故障检测时间缩短 80\% 以上
  \item 在实时业务数据上部署并监控分类与时间序列模型
  \item 实现流式传感器数据的自动化校验与特征提取
\end{itemize}
\cnentry{Sample Tech（远程）}{2021.11 - 2023.02}{机器学习实习生}
\begin{itemize}
  \item 针对技术文档搭建检索系统原型，在内部测试中提升回答准确率
  \item 开发评估工具，持续追踪多个版本间的模型表现
\end{itemize}

\cnsection{教育背景}
\cnentry{Example University}{2021 - 2022}{数据科学硕士}
\begin{itemize}
  \item 毕业项目：面向能源系统的异常检测（统计建模 + 机器学习）
  \item 相关课程：统计机器学习、高级数据库、云计算
\end{itemize}

\cnsection{专业技能}
\begin{itemize}
  \item 编程与数据：Python、SQL、R；Pandas、NumPy
  \item 机器学习与 AI：Scikit-learn、PyTorch；Transformers、LangChain、OpenAI API
  \item 云与部署：Azure、GCP、Docker、Git、GitHub Actions
\end{itemize}

\cnsection{荣誉奖项}
\begin{itemize}
  \item Example AI 黑客松 2025 一等奖
  \item Example University 院长嘉许名单（Dean's List）
\end{itemize}

\end{document}
